\chapter[Aplicação das avaliações]{Aplicação das avaliações} \label{capitulo11}

A aplicação das avaliações constitui uma das principais funcionalidades do sistema, sendo responsável por relacionar um instrumento avaliativo a um determinado aprendente e registrar as informações produzidas durante sua utilização. Essa funcionalidade foi desenvolvida de forma integrada ao gerenciamento dos instrumentos, permitindo que uma mesma avaliação seja utilizada em diferentes atendimentos e que cada aplicação mantenha seu próprio conjunto de respostas e informações contextuais.

No modelo de dados, esse processo é representado pela entidade \textit{aplicacoes}, que estabelece relacionamentos com as entidades \textit{aprendentes}, \textit{avaliacoes} e \textit{users}. Dessa forma, cada aplicação identifica o aprendente avaliado, o instrumento utilizado e o profissional responsável por sua aplicação. Também são armazenadas informações como a data da aplicação, observações, situação do preenchimento e data de finalização.

Para representar o andamento de uma aplicação, foram definidos os estados \textit{PENDING}, \textit{IN\_PROGRESS}, \textit{DRAFT} e \textit{COMPLETED}. Essa estrutura permite diferenciar uma aplicação que ainda não foi iniciada, uma aplicação em andamento, um registro mantido como rascunho e uma aplicação finalizada. O controle desses estados possibilita que o profissional interrompa o preenchimento sem perder os dados já registrados e posteriormente retome o procedimento.

As respostas fornecidas durante a aplicação são armazenadas separadamente na entidade \textit{respostas}. Cada registro associa uma resposta a uma aplicação e a uma questão específica, por meio dos identificadores correspondentes. O campo \textit{valor}, estruturado em formato JSON, permite armazenar diferentes tipos de resposta de acordo com o tipo de questão definido no instrumento. Além disso, foi estabelecida uma restrição de unicidade para a combinação entre \textit{aplicacaoId} e \textit{questionId}, garantindo que uma mesma questão possua apenas uma resposta dentro de determinada aplicação.

A estrutura das questões, definida na entidade \textit{questions}, permite que as aplicações trabalhem com diferentes formatos de resposta. Entre os tipos implementados estão texto, texto longo, valor numérico, data, seleção única, seleção múltipla, escala e respostas do tipo sim/não. Para questões que utilizam opções ou escalas, também podem ser armazenadas as alternativas, os limites e os respectivos rótulos, permitindo que o formulário de aplicação seja construído de acordo com a configuração do instrumento.

Durante o processo de aplicação, o sistema mantém também informações relacionadas ao atendimento, como observações e datas. Esses dados complementam as respostas obtidas no instrumento e contribuem para a contextualização do registro. Ao final do preenchimento, a aplicação pode ser marcada como concluída, mantendo a relação entre o aprendente, o instrumento, o profissional responsável e as respostas registradas.

A separação entre o cadastro do instrumento e suas aplicações permite preservar o histórico dos procedimentos realizados. Assim, o instrumento funciona como uma estrutura reutilizável, enquanto cada aplicação representa uma ocorrência específica de sua utilização. Essa organização é importante para o acompanhamento longitudinal do aprendente, pois diferentes aplicações podem ser associadas ao mesmo indivíduo em momentos distintos, mantendo seus respectivos registros.

A funcionalidade também se relaciona ao processo de elaboração dos relatórios. Uma aplicação concluída pode possuir um relatório associado, no qual o profissional registra sua conclusão e observações. Dessa maneira, os dados obtidos por meio das respostas podem ser posteriormente considerados pelo profissional na elaboração da devolutiva, sem que o sistema realize automaticamente interpretações ou conclusões psicopedagógicas.

Portanto, a aplicação das avaliações foi estruturada para oferecer um fluxo de registro organizado e rastreável, contemplando a seleção do instrumento, sua associação ao aprendente e ao profissional, o preenchimento das questões, o armazenamento das respostas e o controle do estado da aplicação. Com isso, o sistema fornece suporte à documentação das diferentes etapas do processo avaliativo, mantendo a interpretação dos resultados sob responsabilidade do profissional.